\section{Module \texttt{toggle}}

Le module \texttt{toggle} gère les contrôles binaires: case à cocher et interrupteur.

\subsection{Macros de Configuration}

\begin{docbox}{Macros \texttt{CHECKBOX\_CONFIG\_DEFAULT} et \texttt{SWITCH\_CONFIG\_DEFAULT}}
\begin{lstlisting}[language=C]
#define CHECKBOX_CONFIG_DEFAULT ((checkbox_config){ \
    .label = NULL, \
    .active = false, \
    .inconsistent = false, \
    .style = WIDGET_STYLE_CONFIG_DEFAULT, \
    .on_toggled = NULL, \
    .user_data = NULL, \
})

#define SWITCH_CONFIG_DEFAULT ((switch_config){ \
    .label = NULL, \
    .active = false, \
    .state = false, \
    .style = WIDGET_STYLE_CONFIG_DEFAULT, \
    .on_active_changed = NULL, \
    .user_data = NULL, \
})
\end{lstlisting}
\end{docbox}

Ces macros assurent une initialisation contrôlée des états booléens et des callbacks. En laissant \texttt{NULL} par défaut, on évite d'appeler des fonctions non initialisées. Cela limite les crashs lors des changements d'état.

Dans une application interactive, cette discipline d'initialisation est critique: un état binaire mal initialisé peut produire une interface incohérente et des bugs difficiles à diagnostiquer. Les macros réduisent ce risque.

\subsection{Fonctions API}

\begin{lstlisting}[language=C]
GtkWidget *create_switch_row(const switch_config *config, GtkWidget **out_switch) {
    GtkWidget *row = gtk_box_new(GTK_ORIENTATION_HORIZONTAL, 10);
    GtkWidget *label = gtk_label_new(config && config->label ? config->label : "");
    GtkWidget *sw = gtk_switch_new();

    gtk_box_append(GTK_BOX(row), label);
    gtk_box_append(GTK_BOX(row), sw);

    if (config != NULL) {
        gtk_switch_set_active(GTK_SWITCH(sw), config->active);
        gtk_switch_set_state(GTK_SWITCH(sw), config->state);
        apply_widget_style(row, &config->style);

        if (config->on_active_changed != NULL) {
            g_signal_connect(sw, "notify::active", G_CALLBACK(config->on_active_changed), config->user_data);
        }
    }

    if (out_switch != NULL) {
        *out_switch = sw;
    }
    return row;
}
\end{lstlisting}

Le module fournit des composants simples à intégrer tout en conservant un raccordement événementiel explicite.

\texttt{create\_switch\_row} encapsule la composition visuelle (label + switch) et la logique de notification. Le résultat est un composant réutilisable prêt à intégrer dans des formulaires de préférences.

\newpage
